% !PRINTED 458
\weberpage{458}
For the points $P_6$ this is immediately evident, if one sets
[according to \S.\,119, (1)] $x_1 = x_2 = x_3 = 1$, whereby
$f(x_1, x_2, x_3) = 3$, and hence does not vanish. To demonstrate the same for
the $P_8$, one sets, according to \S.\,118 (10):
\[
x_1 = \alpha\gamma\varepsilon^4, \quad x_2 = \beta\gamma\varepsilon,
\quad x_3 = \alpha\beta\varepsilon^2,
\]
whereby, since $\alpha\beta\gamma = \frac{1}{7}$ [\S.\,115, (14)],
\[
f(x_1, x_2, x_3) = \frac{1}{7}(\alpha^3\gamma^2 + \beta^3\alpha^2 + \gamma^3\beta^2)
\]
If one sets in this [according to \S.\,115, (8)]:
\[
\alpha^3 = \alpha\beta\gamma - \alpha^2, \quad \beta^3 = \alpha\beta\gamma - \beta^2,
\quad \gamma^3 = \alpha\beta\gamma - \gamma^2,
\]
one obtains
\weq{6}{f(x_1, x_2, x_3) = \tfrac{1}{7}\alpha\beta\gamma(\alpha^2+\beta^2+\gamma^2) - \tfrac{1}{7}(\alpha^2\gamma^2+\beta^2\alpha^2+\gamma^2\beta^2).}
But from \S.\,115, (6), (7) it follows that
\[
\alpha^2 + \beta^2 + \gamma^2 = 1,
\]
\[
\alpha^2\gamma^2 + \beta^2\alpha^2 + \gamma^2\beta^2 = -\alpha\beta\gamma(\alpha+\beta+\gamma) = \frac{1}{7},
\]
and therefore
\[
f(x_1, x_2, x_3) = \frac{1}{42}
\]
is different from zero.

The intersection points of the 21 principal axes with one another are, in their entirety,
the poles $P_6$ and $P_8$, and consequently two principal axes never
intersect on the fundamental curve. The intersection points of the principal axes with the
fundamental curve can therefore only be four-fold or two-fold poles. To
decide this question, we make the following consideration. If among
the intersection points of the principal axis $A$ with the fundamental curve a $P_2$
occurs, then all four intersection points are distinct from one another and
are all two-fold poles. But if a $P_4$ lies on $A$ and $f$,
then the second $P_4$ lying on $A$ also lies on $f$, and there can be
no other intersection point of $A$ and $f$, because such a point would be
a $P_2$ and would entail four further $P_2$. The four
intersection points of $A$ and $f$ must therefore coincide in pairs, and
$A$ is a double tangent. But now we have seen that the 28
lines $T$ are double tangents, and a curve of the fourth order can
have no more than 28 double tangents. Hence the assumption is
inadmissible, and it follows that the intersection points of the fundamental curve with
the principal axes are two-fold poles.

If, then, we denote the special two-fold poles that lie on the
curve $f$ by $\bar P_2$, we have the following distinguished
systems of points:
%CONT>>
